\section{Introduction}
The starting point for Magda was an earlier project by one of the authors, which used the Gemini API to extract structured offers from weekly supermarket flyers. The estimated API cost in that prototype was approximately EUR~0.08 per flyer. Although this was inexpensive, the prototype required an external model call for every new document. For the Information Extraction seminar, our research question was whether a locally deployed pipeline could match the extraction quality of API-based approaches while eliminating recurring API calls. We assess this through precision, recall and F1.

The original prototype sent images of the flyer pages to Gemini and received structured offers back. For Magda, we divided this process into separate steps. PENNY flyers contain an embedded PDF text layer, so PyMuPDF can extract words and their positions without an OCR step. During fine-tuning, the recognition models learn to label each extracted word as part of a product name, price, quantity or another offer field, or as outside these categories. GBERT and XLM-R use only the word sequence. LiLT also uses the words' positions on the page, while LayoutXLM additionally receives the page image.

Recognising these fields is only part of the task. A flyer page contains several offers, and a price must be assigned to the correct product. Some offers also contain multiple variants or different prices depending on purchase conditions. Magda therefore includes a separate grouping step after recognition, combining the identified entities into offer records. This separation allows us to examine whether an error comes from recognising a field or assigning it to an offer.

Manual annotation was our first approach, but it proved too time-consuming to prepare enough training data within the project timeframe. To build a larger dataset, we used Claude Sonnet 5, a vision-capable large language model, to annotate the flyer pages. The resulting dataset contains 666 pages with 33,457 entity spans. The recognition models were fine-tuned on its training split and evaluated against a separate human-annotated reference.

Following our project proposal, we aimed to build a complete local pipeline that converts flyer PDFs into structured product offers, including names, prices and quantities, using models fine-tuned on our Claude-labelled dataset. To test whether models using layout information improve entity extraction, the evaluation measures precision, recall and F1, compares them with text-only baselines and examines common extraction errors. We also compare the complete pipeline with direct extraction by vision-language models and discuss the trade-offs in extraction quality, cost and local versus API-based processing. Our grouping experiments extend these objectives by examining how recognised entities can be assembled into offer records. Appendix~\ref{app:proposal} summarises how the implementation relates to the proposal.
